% GENERATED FILE. Edit canonical lessons, then run npm run generate:books.
% canonical-source-hash: fnv1a64:7d8eb8919ec5a9ed
% canonical-lessons: PA-C124-amrud, PA-C124-mugphali, PA-C124-sag, PA-C124-lassi, PA-C124-khir

\chapter{Guava, Peanuts, Cooked greens, Lassi, Rice pudding}
\label{ch:pa-guava-peanuts-cooked-greens-lassi-rice-pudding}
\hlchaptermodality{124}

\begin{chapteropening}
\textbf{By the end of this chapter:} \emph{I can say a guava, peanuts, cooked greens, lassi and rice pudding.}

Complete the last lesson of chapter 124: i can say a guava, peanuts, cooked greens, lassi and rice pudding.
\end{chapteropening}

\section[\texorpdfstring{\pa{ਅਮਰੂਦ} (amrūd)}{ਅਮਰੂਦ (amrūd)}]{\pa{ਅਮਰੂਦ} (amrūd) — a guava}

\label{lesson:PA-C124-amrud}



\begin{quote}
Before the new one: say the Punjabi for a watermelon, then the Punjabi for a papaya.
\end{quote}

\subsection*{You'll want to know: \pa{ਅਮਰੂਦ}}
\textbf{\pa{ਅਮਰੂਦ}} — \emph{amrūd} — \textquotedblleft{}a guava\textquotedblright{}.

\begin{grammarlens}[title={What you've built}]
One more word for this chapter.
\end{grammarlens}

\subsection*{Your turn}
\begin{itemize}
\raggedright
  \item \emph{Say it:} \emph{amrūd}
  \item \emph{Say it:} \emph{amrūd}, once more
  \item \emph{Say it:} say \emph{amrūd}
  \item \emph{Recall:} say the Punjabi for a watermelon, then the Punjabi for a papaya, then say \emph{amrūd} again
\end{itemize}

\begin{tcolorbox}[breakable,colback=teal!4,colframe=teal!35!black,title={Before you move on}]
What does \pa{ਅਮਰੂਦ} mean? (\textquotedblleft{}A guava\textquotedblright{}.) Say it once more.
\end{tcolorbox}
\section[\texorpdfstring{\pa{ਮੂੰਗਫਲੀ} (mū̃gphalī)}{ਮੂੰਗਫਲੀ (mū̃gphalī)}]{\pa{ਮੂੰਗਫਲੀ} (mū̃gphalī) — peanuts}

\label{lesson:PA-C124-mugphali}



\begin{quote}
Before the new one: say the Punjabi for a papaya, then the Punjabi for a guava.
\end{quote}

\subsection*{You'll want to know: \pa{ਮੂੰਗਫਲੀ}}
\textbf{\pa{ਮੂੰਗਫਲੀ}} — \emph{mū̃gphalī} — \textquotedblleft{}peanuts\textquotedblright{}.

\begin{grammarlens}[title={What you've built}]
One more word for this chapter.
\end{grammarlens}

\subsection*{Your turn}
\begin{itemize}
\raggedright
  \item \emph{Say it:} \emph{mū̃gphalī}
  \item \emph{Say it:} \emph{mū̃gphalī}, once more
  \item \emph{Say it:} say \emph{mū̃gphalī}
  \item \emph{Recall:} say the Punjabi for a papaya, then the Punjabi for a guava, then say \emph{mū̃gphalī} again
\end{itemize}

\begin{tcolorbox}[breakable,colback=teal!4,colframe=teal!35!black,title={Before you move on}]
What does \pa{ਮੂੰਗਫਲੀ} mean? (\textquotedblleft{}Peanuts\textquotedblright{}.) Say it once more.
\end{tcolorbox}
\section[\texorpdfstring{\pa{ਸਾਗ} (sāg)}{ਸਾਗ (sāg)}]{\pa{ਸਾਗ} (sāg) — cooked greens, saag}

\label{lesson:PA-C124-sag}



\begin{quote}
Before the new one: say the Punjabi for a guava, then the Punjabi for peanuts.
\end{quote}

\subsection*{You'll want to know: \pa{ਸਾਗ}}
\textbf{\pa{ਸਾਗ}} — \emph{sāg} — \textquotedblleft{}cooked greens, saag\textquotedblright{}.

\begin{grammarlens}[title={What you've built}]
One more word for this chapter.
\end{grammarlens}

\subsection*{Your turn}
\begin{itemize}
\raggedright
  \item \emph{Say it:} \emph{sāg}
  \item \emph{Say it:} \emph{sāg}, once more
  \item \emph{Say it:} say \emph{sāg}
  \item \emph{Recall:} say the Punjabi for a guava, then the Punjabi for peanuts, then say \emph{sāg} again
\end{itemize}

\begin{tcolorbox}[breakable,colback=teal!4,colframe=teal!35!black,title={Before you move on}]
What does \pa{ਸਾਗ} mean? (\textquotedblleft{}Cooked greens, saag\textquotedblright{}.) Say it once more.
\end{tcolorbox}
\section[\texorpdfstring{\pa{ਲੱਸੀ} (lassī)}{ਲੱਸੀ (lassī)}]{\pa{ਲੱਸੀ} (lassī) — lassi}

\label{lesson:PA-C124-lassi}



\begin{quote}
Before the new one: say the Punjabi for peanuts, then the Punjabi for cooked greens.
\end{quote}

\subsection*{You'll want to know: \pa{ਲੱਸੀ}}
\textbf{\pa{ਲੱਸੀ}} — \emph{lassī} — \textquotedblleft{}lassi\textquotedblright{}.

\begin{grammarlens}[title={What you've built}]
One more word for this chapter.
\end{grammarlens}

\subsection*{Your turn}
\begin{itemize}
\raggedright
  \item \emph{Say it:} \emph{lassī}
  \item \emph{Say it:} \emph{lassī}, once more
  \item \emph{Say it:} say \emph{lassī}
  \item \emph{Recall:} say the Punjabi for peanuts, then the Punjabi for cooked greens, then say \emph{lassī} again
\end{itemize}

\begin{tcolorbox}[breakable,colback=teal!4,colframe=teal!35!black,title={Before you move on}]
What does \pa{ਲੱਸੀ} mean? (\textquotedblleft{}Lassi\textquotedblright{}.) Say it once more.
\end{tcolorbox}
\section[\texorpdfstring{\pa{ਖੀਰ} (khīr)}{ਖੀਰ (khīr)}]{\pa{ਖੀਰ} (khīr) — rice pudding}

\label{lesson:PA-C124-khir}



\begin{quote}
Before the new one: say the Punjabi for cooked greens, then the Punjabi for lassi.
\end{quote}

\subsection*{You'll want to know: \pa{ਖੀਰ}}
\textbf{\pa{ਖੀਰ}} — \emph{khīr} — \textquotedblleft{}rice pudding\textquotedblright{}.

\begin{grammarlens}[title={What you've built}]
One more word for this chapter.
\end{grammarlens}

\subsection*{Your turn}
\begin{itemize}
\raggedright
  \item \emph{Say it:} \emph{khīr}
  \item \emph{Say it:} \emph{khīr}, once more
  \item \emph{Say it:} say \emph{khīr}
  \item \emph{Recall:} say the Punjabi for cooked greens, then the Punjabi for lassi, then say \emph{khīr} again
\end{itemize}

\begin{tcolorbox}[breakable,colback=teal!4,colframe=teal!35!black,title={Before you move on}]
What does \pa{ਖੀਰ} mean? (\textquotedblleft{}Rice pudding\textquotedblright{}.) Say it once more.
\end{tcolorbox}
